\subsection{Sources and provenance}
GovAlloc-COFOG~v1.0 is assembled exclusively from public, machine-readable official sources that were downloaded through their APIs (scripts in \texttt{pipeline/}; raw files kept verbatim in \texttt{data/raw/}):
\begin{itemize}
\item \textbf{World Development Indicators (WDI)}~\cite{wdi}: \NWdiIndicators{} indicators (GDP, population, inflation, unemployment, tax revenue, central-government debt, health and demographic indicators, infrastructure access, emissions, sector value added), plus four official sector-spending series: education, domestic general-government health, military and R\&D expenditure, each as share of GDP and (for the first three) of government expenditure.
\item \textbf{Worldwide Governance Indicators (WGI)}~\cite{kaufmann2011}: government effectiveness, control of corruption, rule of law and regulatory quality.
\item \textbf{Eurostat \texttt{gov\_10a\_exp}}~\cite{eurostat_cofog} and \textbf{OECD National Accounts Table~11}~\cite{oecd_table11}: general-government expenditure by function (COFOG~\cite{un2000cofog}), which provides the \emph{allocation targets}. OECD data are used only for economies absent from Eurostat (\NOecdAdded{} additional economies, \NOecdRowsComplete{} complete country-years); on the \NOverlapRows{}-row overlap the two sources agree to a median of \NXsrcMedDiff{}~percentage points (Pearson $r\ge \NXsrcMinR{}$ for all functions except communication, Section~\ref{sec:results}).
\item \textbf{IMF GFS expenditure by function}~\cite{imf_gfs}: general-government COFOG divisions (percent of total expenditure) for non-EU/OECD economies, used for the wide panel A10 and cross-checked against Eurostat (median difference $\le0.07$ pp on 723 overlapping country-years).
\end{itemize}
\figCoverage
\subsection{Structure of the released dataset}
The master table (\texttt{govallocxai\_v2\_master.csv}) has \NMasterRows{} rows (\NMasterCountries{} economies $\times$ 2000--2024, plus 25 legacy-only rows for Taiwan) and \NMasterCols{} columns: identifiers and country groupings, \NFeatNum{} features with a per-cell imputation flag, target series with source labels, and \emph{all columns of the legacy v1 dataset} (prefix \texttt{legacy\_}; country and year serve as keys) so that the earlier dataset is preserved and directly comparable. Two modelling panels are derived from it:
\begin{description}
\item[Panel A (COFOG; ``A16'' in the results).] \NCofogRows{} country-years from \NCofogCountries{} economies with a \emph{complete} 16-function breakdown. Targets are the share of total government expenditure (\%) of 16 functions: general public services, defence, public order and safety, agriculture, energy, industry and construction, transport, communication, R\&D (economic affairs), environment, housing and community amenities, water supply, health, recreation and culture, education and social protection, plus an explicit \emph{other/unallocated} residual so that the vector lies on the simplex (sums to 100\%). The 16 ``functions'' mix COFOG levels: general public services, defence, public order, environment, health, recreation, education and social protection are COFOG divisions; agriculture, energy, industry and construction (04.4 and 04.7), transport, communication and R\&D are groups of division 04 (economic affairs); water supply is group 06.3, and housing and community amenities is division 06 \emph{excluding} 06.3, so nothing is counted twice. The residual collects the remaining economic-affairs groups (general economic affairs 04.1 and 04.9 n.e.c.); it averages \ResidualMean\% of expenditure (median \ResidualMedian\%, maximum \ResidualMax\%) and is predicted and scored like a function.
\item[Panel A10 (wide COFOG).] \NTenRows{} country-years from \NTenEcon{} economies with a complete ten-division COFOG breakdown from Eurostat, OECD and IMF (priority in that order), including 35 middle- and low-income economies under the World Bank's \emph{current} classification, which is applied to all years. TVD values are not comparable between A16 (17 components) and A10 (10 components).
\item[Panel B (WDI).] \NWdiRows{} country-years from \NWdiCountries{} economies with at least one of the four official sector series (global coverage including low-income countries).
\end{description}
\subsection{Missing-value policy}
Targets are \emph{never} imputed. Features are filled in three ordered, flagged steps: (i) linear interpolation within a country series (maximum three consecutive years), (ii) edge fill (maximum three years), (iii) median of the region $\times$ income group $\times$ year cell. These steps are applied once to the whole table \emph{before} any cross-validation split, so interpolation can use later values and the cell medians pool all economies, including held-out ones; imputation is transductive and only features, never targets, are affected (Section~\ref{sec:disc}). Each imputed cell carries an \texttt{imp\_*} flag and each row an imputed-feature count. Features with high structural missingness (poverty headcount, Gini, adult literacy: 65--81\,\% of cells imputed) are \emph{excluded} from all models; \NFeatModel{} features remain (Table~\ref{tab:features}).
\subsection{Legacy dataset (v1)}
The v1 dataset (775 rows, 31 economies, 59 columns) is preserved verbatim. Its \emph{features} match the World Bank API values (exactly for GDP, population, life expectancy, growth, unemployment and internet use, and to within 0.4\% on 90\% of rows for inflation, on \NLegacyMatched{} matched rows), whereas its fifteen sector columns are mostly synthetic (Section~\ref{sec:audit}). v1 is therefore used only as an audited negative control, never as training evidence for the main results.
